% Page 027

\seite{27}

\jahresueberschrift{Das Schuljahr 1899}
\pstart
Aufgenommen wurden sieben Neulinge, somit\ob
die Schule in Gesamtstand Schülerzahl von 50\ob
Kindern, 21 Mädchen und 29 Knaben.\ob
Besorgt wird der Schulkinder, wurde die\ob
Schule zu Einrichtung, Briefe beim zur Oster\ohb
ferien übernehmen.

\pend
\abschnitt{Rohbeheitung [?]}
\pstart

\margmark{Erlaß 11.9.99          I\\Fleuth}%
n der Lehrer des Nachmittags Seffern zu einer\ob
gewöhnungstätigen Weisung unterliegt. Der Schul\ohb
verbindung mit Lehrer der hiesigen Schule abzulegen.\ob
Über diesen Grundes ist in der hiesigen Schule Halt\ob
begründet, nicht eingerichtet mit dem Maßnahmen, daß ab\ohb
wechselnd der Unterricht das Komitteliges Statt\ob
Aufstellung in der hiesigen Schule abgehalten wird —\ob
da der Lehrer in Seffern (Kohlbrecher) auf 14\ob
Tagen für dienstunfähigerlich erklärt wurde,\ob
konnte er seinen Dienst in Seffern wieder\ob
aufnehmen und somit hatte die Vertretung\ob
ein Ende\unsicher.

Ausfall des Unterrichts wegen Hitze.

Im Monat August und September mußte der\ob
Unterricht eingestellt, wegen zu großer Hitze\ob
eingestellt werden.
\pend
